% !TeX root = ../../Book3.tex
\section{Cauchy 积分定理}
\label{b3:sec:0202}

局部 Taylor 展开给出小圆盘内的原函数。若曲线穿过许多圆盘，不同原函数的端点差能否接起来？若曲线围住一个孔洞，它们又为什么可能留下非零总和？本节先把积分沿连续路径定义清楚，再用参数方格说明连续变形怎样保持积分，最后把路径独立性转回原函数的存在。

\subsection{三角形积分与局部原函数}
\label{b3:subsec:020201}

若 \(f\) 在包含一个闭三角形 \(T\) 的开集上全纯，\cref{b3:lem:01-goursat-local}已经证明 \(\int_{\partial T}f(z)\,\mathrm{d}z=0\)。现在全纯即解析已经建立，可以直接逐项积分，得到小圆盘上的原函数。

\Needspace{6\baselineskip}
\begin{lemma}[局部原函数]\label{b3:lem:02-local-primitive}
设 \(a\in\mathbb C\)、\(r>0\)，\(f:D(a,r)\to\mathbb C\) 全纯。则存在全纯函数 \(F:D(a,r)\to\mathbb C\)，满足 \(F'=f\)；还可要求 \(F(a)=0\)，此时 \(F\) 唯一。
\end{lemma}
\begin{proof}
由\cref{b3:thm:01-holomorphic-analytic}，有
\[
 f(z)=\sum_{n=0}^{\infty}c_n(z-a)^n
 \quad(|z-a|<r).
\]
令
\[
 F(z)=\sum_{n=0}^{\infty}\frac{c_n}{n+1}(z-a)^{n+1}.
\]
对每个 \(0<\rho<r\)，这个级数的绝对值在 \(|z-a|\le\rho\) 上由
\(\rho\sum_{n\ge0}|c_n|\rho^n\) 控制，因而在每个较小闭圆盘上一致绝对收敛。\Cref{b3:thm:01-power-differentiation}允许逐项求导，得到 \(F'=f\)，而 \(F(a)=0\) 直接来自级数。两个这样的原函数之差导数为零；圆盘连通，故其差为常数，再在 \(a\) 处求值即得唯一性。
\end{proof}

在一般开集 \(U\) 内，每个点都有闭包仍包含于 \(U\) 的小圆盘，所以局部原函数总能取得。不过，相邻圆盘上的原函数未必以同一常数归一化。积分只使用端点差，常数在每段内都会消失；这使我们不必先解决全局拼合，就能沿一条路径继续计算。

\subsection{连续路径上的局部积分}
\label{b3:subsec:020202}

连续变形产生的中间路径未必分段光滑，甚至可能没有有限长度。此时不能把不存在的导数放进参数积分。但是只要被积函数在每个小圆盘有原函数，沿一段路径的贡献已经由端点决定。这给出一种不对路径求导的定义。

\ProseParagraphBreak

先说明怎样选出有限分段。紧度量空间 \(K\) 的任意开覆盖都有一个数 \(\delta>0\)，使直径小于 \(\delta\) 的每个非空子集落在一个覆盖成员中。这样的数称为这个覆盖的\term[lebesgue-shu]{Lebesgue 数}[Lebesgue number]。为证明它，对每个 \(x\in K\) 选 \(r_x>0\)，使相对球 \(B_K(x,2r_x)\) 落在某个覆盖成员内；由紧性，从 \(B_K(x,r_x)\) 中取有限子覆盖，并令 \(\delta\) 为所取有限个半径的最小值。若小集合含有一点 \(y\in B_K(x,r_x)\)，则其中任一点到 \(x\) 的距离小于 \(\delta+r_x\le2r_x\)，故整个集合属于同一个覆盖成员。下面分别把 \(K\) 取成参数区间与参数方块。

\Needspace{6\baselineskip}
\begin{lemma}[局部积分与连续路径]\label{b3:lem:02-continuous-path-integral}
设 \(U\subseteq\mathbb C\) 为开集，\(f:U\to\mathbb C\) 全纯，\(\gamma:[0,1]\to U\) 连续。可以选择分割 \(0=t_0<\cdots<t_N=1\) 以及圆盘 \(D_k\subset U\)，使 \(\gamma([t_{k-1},t_k])\subset D_k\)。在每个 \(D_k\) 上选择 \(f\) 的一个原函数 \(F_k\)，令
\[
 I_f(\gamma)\coloneqq\sum_{k=1}^N
 \bigl(F_k(\gamma(t_k))-F_k(\gamma(t_{k-1}))\bigr).
\]
这个数与分割、圆盘和原函数的选择无关。它在严格递增连续双射换参下不变，在路径拼接下可加，在反向时变号；若 \(\gamma\) 分段光滑，则 \(I_f(\gamma)=\int_\gamma f(z)\,\mathrm{d}z\)。
\end{lemma}
\begin{proof}
对紧像集中的每一点选一个包含于 \(U\) 的圆盘。圆盘的逆像组成 \([0,1]\) 的开覆盖，取其 Lebesgue 数 \(\delta>0\)，再选择每段长度小于 \(\delta\) 的有限分割。每段路径的像就完整地落在一个圆盘中；前一引理提供该圆盘上的原函数。

固定一套圆盘和原函数时，加入分点只是把一个端点差拆成两个，所以不改变和。现在比较同一参数段所能使用的两个圆盘 \(D,E\)，并设它们的原函数为 \(F,G\)。整段路径的像包含在 \(D\cap E\) 中。两个圆盘的非空交是凸集，因而连通；\((F-G)'=0\)，所以 \(F-G\) 在交集中为常数。于是
\[
 F(\gamma(t_k))-F(\gamma(t_{k-1}))
 =G(\gamma(t_k))-G(\gamma(t_{k-1})).
\]
对任意两套选择，把两套分割的全部分点合并。细分后的每一小段同时处于原来两套所选的圆盘中，逐段应用这个等式，再相加，就得到相同的总和。

对严格递增连续双射换参，把原分点换成其逆像，所有端点及次序不变，故总和不变。反向使每段的两个端点交换，拼接使两组端点差依次排列，因此分别变号与相加。若路径分段光滑，把其光滑分点加入分割；每段的端点差由\cref{b3:prop:02-primitive-endpoints}等于通常积分，所以两种定义一致。
\end{proof}

以下仍将这个数写成 \(\int_\gamma f(z)\,\mathrm{d}z\)。它针对具有局部原函数的被积函数定义，不是在任意不可求长路径上积分任意连续函数的方法。对这里的全纯函数，每一小段还可替换成连接其端点的直线段：圆盘凸性使线段留在原圆盘内，端点公式保证积分相同。这样得到一条有限折线路径，适合实际计算。

\ProseParagraphBreak

不同圆盘的端点差不能未经核对就全部望远镜相消。以 \(1/z\) 沿单位圆为例，每个短圆弧都可选连续辐角，并用相应对数作原函数；沿逆时针方向依次衔接后，辐角从 \(0\) 增加到 \(2\pi\)，端点差之和为 \(2\pi\mathrm{i}\)。若存在同一个全局原函数，起点与终点相同才会迫使总和为零。下一小节说明何种路径变形可以保证这种总和消失。

\subsection{单连通区域}
\label{b3:subsec:020203}

用连续映射 \(H:[0,1]^2\to U\) 表示路径的连续变形：固定第二个参数 \(t\) 后，\(s\mapsto H(s,t)\) 是变形到时刻 \(t\) 的路径。这种变形称为\term[tonglun]{同伦}[homotopy]。若所有路径具有固定的起点和终点，便有 \(H(0,t)=p\)、\(H(1,t)=q\)；若只要求每条路径闭合，则要求 \(H(0,t)=H(1,t)\)，共同端点可以移动。

\Needspace{6\baselineskip}
\begin{proposition}[同伦下的积分不变性]\label{b3:prop:02-homotopy-invariance}
设 \(U\subseteq\mathbb C\) 为开集，\(f:U\to\mathbb C\) 全纯，\(H:[0,1]^2\to U\) 连续。定义四条路径
\[
 \gamma_j(s)=H(s,j),
 \quad \eta_j(t)=H(j,t)
 \quad(j=0,1).
\]
沿连续路径的积分按\cref{b3:lem:02-continuous-path-integral}定义，则
\[
 \int_{\gamma_0}f(z)\,\mathrm{d}z
 +\int_{\eta_1}f(z)\,\mathrm{d}z
 -\int_{\gamma_1}f(z)\,\mathrm{d}z
 -\int_{\eta_0}f(z)\,\mathrm{d}z=0.
\]
特别地，若 \(H(0,t)\)、\(H(1,t)\) 分别与 \(t\) 无关，或对所有 \(t\) 有 \(H(0,t)=H(1,t)\)，则
\(\int_{\gamma_0}f(z)\,\mathrm{d}z=\int_{\gamma_1}f(z)\,\mathrm{d}z\)。端路径分段光滑时，这些积分等于通常的曲线积分。
\end{proposition}
\begin{proof}
用包含于 \(U\) 的原函数圆盘覆盖紧集 \(H([0,1]^2)\)。其逆像是参数方块的开覆盖，取 Lebesgue 数 \(\delta>0\)。选择整数 \(N\) 使 \(\sqrt2/N<\delta\)，将方块分成边长 \(1/N\) 的 \(N^2\) 个闭小方格。每格的直径小于 \(\delta\)，故整格的像位于一个原函数圆盘内。

取一个小格，按逆时针顺序记其四个顶点为 \(v_0,v_1,v_2,v_3\)，并记 \(v_4=v_0\)。若 \(F\) 是覆盖该格像的原函数，则四条边经 \(H\) 映成连续路径，其积分总和是
\[
 \sum_{k=0}^3\bigl(F(H(v_{k+1}))-F(H(v_k))\bigr)=0.
\]
这个计算不要求 \(H\) 可微或单射，也不要求格子的像围成一个区域。

把所有小格的等式相加。每条内部边恰出现两次，一次从某个顶点走向另一个顶点，一次反向行走；其像是同一条参数路径的两个方向。即使相邻格子使用不同原函数，前一引理的良定义性也保证两次积分互为相反数，所以它们相消。余下外边界的走向依次为：底边从左到右，即 \(\gamma_0\)；右边从下到上，即 \(\eta_1\)；顶边从右到左，即 \(\gamma_1^-\)；左边从上到下，即 \(\eta_0^-\)。这恰好给出所写等式。

定端点时，两条竖边映为常路径，积分都是零。闭路径同伦时，两条竖边映为同一连续路径，故两个积分抵消。两种情形都剩下 \(\gamma_0\) 与 \(\gamma_1\) 的积分相等。
\end{proof}

这里的紧性用于找到一个统一的细分尺度；局部原函数用于计算每格的四个端点差；反向与拼接规则用于消去公共边。无须对 \(H\) 求偏导，也无须把弯曲或重叠的格子像看作可以使用平面面积公式的区域。

\Needspace{6\baselineskip}
\begin{definition}\label{b3:def:02-simply-connected}
区域 \(U\subseteq\mathbb C\) 称为\term[danliantong]{单连通}[simply connected]，是指对每条闭连续路径 \(\gamma:[0,1]\to U\)，都存在点 \(p\in U\) 及连续映射 \(H:[0,1]^2\to U\)，满足
\[
 H(s,0)=\gamma(s),
 \quad H(s,1)=p,
 \quad H(0,t)=H(1,t)
 \quad(0\le s,t\le1).
\]
也就是说，每条闭路径都能在 \(U\) 内经过闭路径连续收缩到一点。
\end{definition}

每个非空凸开集都单连通：固定 \(p\in U\)，取
\(H(s,t)=(1-t)\gamma(s)+tp\)，凸性保证所有变形点仍在 \(U\) 内，闭合关系由 \(\gamma(0)=\gamma(1)\) 保留。圆盘、半平面都属此例。在一般区域，是否存在这样的收缩就是积分定理所需的拓扑条件。

\ProseParagraphBreak

\term[cauchy-jifen-dingli]{Cauchy 积分定理}[Cauchy's integral theorem]现在由同伦不变性直接得到。

\Needspace{6\baselineskip}
\begin{theorem}[Cauchy，单连通区域]\label{b3:thm:02-cauchy-simply-connected}
设 \(U\subseteq\mathbb C\) 是单连通区域，\(f:U\to\mathbb C\) 全纯。则 \(U\) 内每条闭分段光滑曲线 \(\gamma\) 都满足
\[
 \int_\gamma f(z)\,\mathrm{d}z=0.
\]
\end{theorem}
\begin{proof}
按定义，把 \(\gamma\) 在 \(U\) 内通过闭路径连续收缩到常路径。由\cref{b3:prop:02-homotopy-invariance}，其积分等于常路径的积分，后者为零。连续路径积分与 \(\gamma\) 的通常曲线积分一致，所以得到结论。
\end{proof}

\begin{figure}[htbp]
\centering
\begin{tikzpicture}[>=stealth,scale=0.95,font=\small]
 \draw[step=0.75,gray!65] (0,0) grid (3,3);
 \fill[gray!20] (0.75,0.75) rectangle (1.5,1.5);
 \draw (0,0) rectangle (3,3);
 \draw[->,thick] (0.7,0)--(1.5,0);
 \draw[->,thick] (3,0.8)--(3,1.5);
 \draw[->,thick] (2.3,3)--(1.5,3);
 \draw[->,thick] (0,2.2)--(0,1.5);
 \draw[->] (1.44,0.87)--(1.44,1.37);
 \draw[->] (1.56,1.37)--(1.56,0.87);
 \node[below] at (1.5,0) {\(H(s,0)=\gamma(s)\)};
 \node[above] at (1.5,3) {\(H(s,1)=p\)};
 \node[rotate=90,above] at (0,1.5) {\(H(0,t)\)};
 \node[rotate=90,below] at (3,1.5) {\(H(1,t)\)};
 \node at (1.12,1.12) {\(Q\)};
 \draw[->,thick] (3.6,1.6)--(4.7,1.6) node[midway,above] {\(H\)};
 \draw (6,1.6) circle (1.1);
 \draw[fill=gray!20] (5.45,1.35)..controls (5.7,2.05) and (6.45,2.2)..(6.55,1.6)
   ..controls (6.4,1.05) and (5.75,0.95)..cycle;
 \node at (6,1.55) {\(H(Q)\)};
 \node[above] at (6,2.7) {原函数圆盘 \(D_j\)};
\end{tikzpicture}
\caption[同伦方格中的边界相消]{细分的是参数方块，不是 \(U\)。每格的像包含于某个原函数圆盘；右图仅示意包含关系，像不必是简单区域。相邻小格的公共边方向相反，外侧两竖边的像相同而方向相反，顶边的像为一点，最后只剩底边 \(\gamma\)。}
\label{b3:fig:02-homotopy-grid}
\end{figure}

\Cref{b3:fig:02-homotopy-grid}画出收缩同伦的情形。细分发生在参数方块，灰格的像只需位于一个原函数圆盘内；它不必是简单图形。闭合关系使外侧两竖边的积分抵消，常值顶边没有贡献，所以最后底边积分为零。

\ProseParagraphBreak

去心平面不是单连通的，因为若单位圆能在其中收缩，对全纯函数 \(1/z\) 使用同伦不变性，就会得到 \(2\pi\mathrm{i}=0\)。同样，在圆环 \(r<|z|<R\) 内，绕中心的圆周不能收缩到点，但半径介于 \(r\) 和 \(R\) 的任意两个正向圆周可以相互变形：
\[
 H(s,t)=\bigl((1-t)\rho_0+t\rho_1\bigr)e^{2\pi\mathrm{i}s},
 \quad r<\rho_0,\rho_1<R.
\]
所以任意在圆环内全纯的函数沿这两条圆周的积分相同。积分不变并不总意味着积分为零；只有能变形成常路径时，才取得零值。

\subsection{原函数判别与 Morera 定理}
\label{b3:subsec:020204}

若所有闭路积分都为零，从一个固定起点出发积分，就不会因路径选择不同而得到两个值。这给出了寻找全局原函数的自然构造。所需的被积函数条件可以放宽为连续，全纯性会成为构造的结果。

\Needspace{6\baselineskip}
\begin{proposition}[原函数判别准则]\label{b3:prop:02-primitive-criterion}
设 \(U\subseteq\mathbb C\) 为区域，\(f:U\to\mathbb C\) 连续。则 \(f\) 在 \(U\) 上有全纯原函数，当且仅当它沿 \(U\) 内每条闭分段光滑曲线的积分为零。上述条件还等价于：任意两条起点相同、终点相同的分段光滑路径给出相同的积分。
\end{proposition}
\begin{proof}
若有原函数，端点公式立即给出路径独立性，也给出闭路积分为零。路径独立性反过来蕴含闭路积分为零，因为闭路可以与同一起终点的常路径比较。以下由闭路积分消失构造原函数。

固定 \(a\in U\)。先验证每一点都可从 \(a\) 沿有限折线路径到达。记这些可达点组成的集合为 \(E\)，则 \(a\in E\)。对任意 \(z\in U\)，取小圆盘 \(D(z,r)\subset U\)。若圆盘内有一个可达点，就可以从该点续接线段到圆盘内任意点，所以整个圆盘都属于 \(E\)；若没有可达点，整个圆盘属于补集。因此 \(E\) 与 \(U\setminus E\) 都在 \(U\) 中开，由连通性得到 \(E=U\)。

对 \(z\in U\)，选一条由 \(a\) 到 \(z\) 的折线路径，定义 \(F(z)\) 为其上 \(f\) 的积分。两条这样的路径与其中一条的反向拼成闭路，闭路积分为零，所以定义与选择无关。固定 \(z\) 并取 \(D(z,r)\subset U\)，当 \(0<|h|<r\) 时，可在到达 \(z\) 后续接线段到 \(z+h\)，故
\[
 \frac{F(z+h)-F(z)}h
 =\int_0^1f(z+th)\,\mathrm{d}t.
\]
右边与 \(f(z)\) 之差的模不超过
\(\sup_{|w-z|\le|h|}|f(w)-f(z)|\)，由连续性趋于零。因此 \(F'(z)=f(z)\)，\(F\) 在 \(U\) 全纯，构造完成。
\end{proof}

结合 Cauchy 积分定理，单连通区域上的每个全纯函数都有原函数。但“某一个函数有原函数”不能反推区域单连通：例如 \(f=1\) 在任意区域都有原函数 \(F(z)=z\)。去心平面上的 \(1/z^2\) 也有原函数 \(-1/z\)，而 \(1/z\) 没有。孔洞只提供可能的障碍，是否产生非零闭路积分还取决于被积函数。

\ProseParagraphBreak

原函数构造还能局部使用。如果只知道所有小三角形的边界积分为零，便在圆盘内以线段连接固定基点；三角形条件恰好保证改变终点时可以续接一条短线段。这就是\term[morera-dingli]{Morera 定理}[Morera's theorem]的证明方法。

\Needspace{6\baselineskip}
\begin{theorem}[Morera]\label{b3:thm:02-morera}
设 \(U\subseteq\mathbb C\) 为开集，\(f:U\to\mathbb C\) 连续。若对每个闭三角形 \(T\subset U\)，都有
\(\int_{\partial T}f(z)\,\mathrm{d}z=0\)，则 \(f\) 在 \(U\) 上全纯。三角形边界取正向。
\end{theorem}
\begin{proof}
任取 \(a\in U\) 及圆盘 \(D(a,r)\subset U\)，对 \(z\) 在圆盘内定义
\(F(z)=\int_{[a,z]}f(\zeta)\,\mathrm{d}\zeta\)。当 \(z+h\) 仍在圆盘内时，凸性保证以 \(a,z,z+h\) 为顶点的三角形位于圆盘中。它的边界积分为零，因而
\[
 F(z+h)-F(z)=\int_{[z,z+h]}f(\zeta)\,\mathrm{d}\zeta.
\]
三点共线时，同一式子直接由有向线段积分的可加性得到。于是对 \(h\ne0\)，差商仍为 \(\int_0^1f(z+th)\,\mathrm{d}t\)，由连续性趋于 \(f(z)\)。所以 \(F\) 全纯且 \(F'=f\)。由第一章已经证明的全纯即解析，\(F\) 的导数也全纯，从而 \(f\) 在这个圆盘上全纯。每一点都能选择这样的圆盘，故结论在 \(U\) 上成立。
\end{proof}

连续性在最后的差商极限中不可省略。例如在复平面上令 \(f(0)=1\)，其余点取 \(f(z)=0\)。任意非退化三角形的各条边最多在一个参数点遇到原点，改变有限个参数点的值不改变实积分，所以所有三角形边界积分仍为零；但 \(f\) 在原点不连续，更不全纯。积分条件需要连续性才能恢复原来的点值。

\ProseParagraphBreak

接下来把上下两侧的全纯函数跨过一条直线粘合。这里不能对交界处的未知导数直接使用 Cauchy--Riemann 方程；改用 Morera，只须证明切割后的边界积分能在极限中相接。

\Needspace{6\baselineskip}
\begin{proposition}[跨直线粘合]\label{b3:prop:02-line-gluing}
设 \(U\subseteq\mathbb C\) 为开集，\(f:U\to\mathbb C\) 连续，并在 \(U\setminus\mathbb R\) 上全纯。则 \(f\) 在 \(U\) 上全纯。
\end{proposition}
\begin{proof}
由 Morera 定理，只需证明任意闭三角形 \(T\subset U\) 的边界积分为零。先记录连续性对短距离移动线段的作用。设 \(a_m\to a\)、\(b_m\to b\)，且全部线段都位于紧集 \(T\) 内，取线性参数化
\(\sigma_m(t)=a_m+t(b_m-a_m)\)、\(0\le t\le1\)。它一致趋于 \(\sigma(t)=a+t(b-a)\)，导数常值也趋于 \(b-a\)。由于 \(f\) 在 \(T\) 上一致连续且有界，
\[
 \int_{[a_m,b_m]}f(z)\,\mathrm{d}z
 =(b_m-a_m)\int_0^1 f(\sigma_m(t))\,\mathrm{d}t
 \longrightarrow\int_{[a,b]}f(z)\,\mathrm{d}z.
\]
若 \(a=b\)，右端按长度为零的常线段解释，结论仍成立。

令 \(T_+=T\cap\{\operatorname{Im}z\ge0\}\)，先假定它有非空内部。对充分小的 \(\varepsilon>0\)，令
\[
 P_\varepsilon=T\cap\{\operatorname{Im}z\ge\varepsilon\}.
\]
这是一个非退化凸多边形，闭包包含于 \(U\setminus\mathbb R\)。从其中一个顶点向不相邻顶点连线，可将它剖成有限个三角形；\cref{b3:lem:01-goursat-local}使每个三角形的边界积分为零，内部公共边相消，得到 \(\int_{\partial P_\varepsilon}f(z)\,\mathrm{d}z=0\)。

当 \(\varepsilon\downarrow0\) 时，\(P_\varepsilon\) 的各条边由原三角形边的截取部分及一条水平截线组成。各端点沿固定线段连续趋向 \(T_+\) 的相应边端点；某条边若缩成一点，其积分趋于零。新增水平截线趋于 \(T_+\) 在实轴上的底边，或者趋于一个顶点。对这有限条线段分别应用刚证的端点连续性，并保留逆时针方向，得到
\(\int_{\partial T_+}f(z)\,\mathrm{d}z=0\)。若 \(T_+\) 完全位于上半平面，则充分小时截取不再改变它，同一结论更直接。

对 \(T_-=T\cap\{\operatorname{Im}z\le0\}\)，用水平线 \(\operatorname{Im}z=-\varepsilon\) 从下侧截取，得到同样的零积分结论。空的部分不作贡献；若某一部分退化成线段或点，其正反两向积分相消。最后将上下两部分的边界积分相加，公共实线段分别向左、向右行走，因而完全抵消；其余边正好组成 \(\partial T\)。于是 \(\int_{\partial T}f(z)\,\mathrm{d}z=0\)，Morera 定理给出全纯性。
\end{proof}

交界处函数值的连续性使两侧的积分贡献相消。全纯性成立以后，所有导数的相容性才由解析性得到。若上侧取常数 \(1\)、下侧取常数 \(0\)，两侧都全纯，却无法连续粘合，恰好缺少命题的假设。对任意其他直线，先作平移和非零复数旋转把它变成实轴，再应用同一结论即可；后面的反射原理将用到这种局部粘合。
